\chapter[A Primer on Statistical Mechanics*]{A Primer on Statistical Mechanics\texorpdfstring{$^*$}{*}}
\label{app:statmech}
% ============================================================================
% APPENDIX WRITTEN FROM SCRATCH (asterisk convention): requested by the author on
% 2026-07-06 and written on 2026-07-07 from standard references (Sethna; Gardiner),
% deliberately using the NOTATION OF THE COURSE NOTES: β = (k_B T)^{-1}, the
% Boltzmann formula ρ(E) = N e^{-βE}/Z as in Chapter 2 (Einstein's second paper),
% ensemble averages ⟨·⟩, and the OU parameters (κ, D) of Chapter 4.
% Content: ensembles and the Boltzmann-Gibbs distribution; partition function and
% free energy; equipartition; barometric and Maxwell-Boltzmann distributions;
% stationary Fokker-Planck states, detailed balance, and the Einstein relation.
% ============================================================================

The main text borrows, at several key moments, results from equilibrium statistical mechanics: the Boltzmann formula in Einstein's second paper (Chapter~2), the equipartition theorem in Langevin's derivation (Chapter~2), the Maxwell--Boltzmann distribution as the stationary state of the Ornstein--Uhlenbeck process (Chapter~4), and the notions of equilibrium and detailed balance behind the wealth distribution (Chapter~8).
This appendix collects those results in one place, at primer level: enough to make the book self-contained, not a substitute for a statistical mechanics course\cite{Sethna2021, Gardiner2009}.
The notation follows the one used in the main text.


\section{Statistical ensembles and the Boltzmann distribution}
\label{app:ensembles}

A macroscopic system has an astronomical number of microscopic degrees of freedom, and its exact microstate is neither knowable nor interesting.
Statistical mechanics replaces the single system by an \emph{ensemble}: a mental collection of very many copies, identical at the macroscopic level and distributed over the possible microstates, so that macroscopic quantities become ensemble averages $\langle\,\cdot\,\rangle$ --- exactly the averages over realizations that we take throughout the book for stochastic processes (Chapter~3).

For an \emph{isolated} system, the fundamental postulate is that all microstates compatible with the fixed total energy are equally probable (the \emph{microcanonical} ensemble).
The workhorse of this book, however, is the \emph{canonical} ensemble: a system in contact with a much larger environment (a \emph{bath}) at temperature $T$, with which it can exchange energy.
The probability of finding the system in a microstate of energy $E$ then follows from the postulate applied to the total (system + bath): it is proportional to the number of bath microstates left over, and expanding the bath entropy to first order in the small energy $E$ ceded to the system gives the celebrated exponential weight
\begin{equation}\label{eq:boltzmann-weight}
P(\text{microstate with energy } E) \;\propto\; e^{-\beta E},
\qquad
\beta = \frac{1}{k_B T},
\end{equation}
the \emph{Boltzmann distribution}, with $k_B \simeq 1.38\times10^{-23}\,\mathrm{J/K}$ the Boltzmann constant.
High-energy states are exponentially penalized, the more severely the lower the temperature; at high temperature the weight flattens and all states become nearly equiprobable.

Normalizing the weight defines the \emph{partition function} $\mathcal{Z}$:
\begin{equation}\label{eq:partition-function}
P(E) = \frac{e^{-\beta E}}{\mathcal{Z}},
\qquad
\mathcal{Z} = \sum_{\text{microstates}} e^{-\beta E}
\quad\text{or}\quad
\mathcal{Z} = \int d(\text{states})\, e^{-\beta E}
\end{equation}
in the discrete and continuous cases respectively.
In the main text the same statement appears in the number-density form used by the notes: for an ensemble of $N$ identical particles,
\begin{equation}\label{eq:boltzmann-density}
\rho(E) = N\,\frac{e^{-\beta E}}{\mathcal{Z}},
\end{equation}
which is Equation~\eqref{eq:boltzmann-weight} multiplied by the total number of particles --- the form invoked in Einstein's second paper (Chapter~2).

The partition function is far more than a normalization.
Averages of the energy follow from its derivatives,
\begin{equation}\label{eq:z-derivatives}
\langle E\rangle = -\frac{\partial}{\partial\beta}\ln\mathcal{Z},
\qquad
\mathrm{Var}[E] = \frac{\partial^2}{\partial\beta^2}\ln\mathcal{Z},
\end{equation}
so $\ln\mathcal{Z}$ generates the energy statistics exactly as the logarithm of the characteristic function generates the moments of a random variable (Chapter~2) --- the two formalisms are the same mathematics.
The quantity
\begin{equation}\label{eq:free-energy}
F = -k_B T\,\ln\mathcal{Z}
\end{equation}
is the \emph{free energy}, the fundamental thermodynamic potential of the canonical ensemble: the system in equilibrium minimizes not its energy but $F = \langle E\rangle - TS$, the balance of energy against entropy that drives every phase transition.

\medskip

Two applications of the Boltzmann weight are used repeatedly in the book.
For a classical particle in an external potential, the energy splits as $E = V(x) + \tfrac{mv^2}{2}$, and the weight factorizes into a position part and a velocity part.
The position part gives
\begin{equation}\label{eq:barometric-app}
\rho(x) \;\propto\; e^{-\beta V(x)},
\end{equation}
which for the gravitational potential $V(x) = mgx$ is the \emph{barometric distribution} $\rho(x) = \rho(0)\,e^{-\beta mgx}$ of Chapter~2: the density of particles thins out exponentially with height, the competition between gravity and thermal agitation.
The velocity part gives
\begin{equation}\label{eq:maxwell-boltzmann}
P^{\mathrm{M.B.}}(v) \;\propto\; e^{-\beta m v^2/2},
\end{equation}
the \emph{Maxwell--Boltzmann distribution}: a Gaussian in the velocity, of zero mean and variance
\begin{equation}\label{eq:mb-variance}
\mathrm{Var}[v] = \frac{k_B T}{m},
\end{equation}
which is the equilibrium state that the Ornstein--Uhlenbeck process of Chapter~4 relaxes to.


\section{The equipartition theorem}
\label{app:equipartition}

The Maxwell--Boltzmann variance~\eqref{eq:mb-variance} is one instance of a remarkably general statement: at equilibrium, \emph{every quadratic degree of freedom carries the same average energy}, namely $\tfrac12 k_B T$.
This is the \emph{equipartition theorem}, and it is the single statistical-mechanics input of Langevin's derivation in Chapter~2.

The proof is one Gaussian integral.
Let the energy depend quadratically on some degree of freedom $q$ (a velocity component, a coordinate in a harmonic potential), $E = a q^2 + E'$, where $a>0$ and $E'$ collects everything else.
By the Boltzmann weight~\eqref{eq:partition-function} the average of the quadratic term is
\begin{equation}\begin{split}\label{eq:equipartition-proof}
\langle a q^2\rangle
&= \frac{\displaystyle\int_{-\infty}^{\infty} dq\; a q^2\, e^{-\beta a q^2}}{\displaystyle\int_{-\infty}^{\infty} dq\; e^{-\beta a q^2}}\\
&= -\frac{\partial}{\partial\beta}\,\ln\!\int_{-\infty}^{\infty} dq\; e^{-\beta a q^2}
= -\frac{\partial}{\partial\beta}\,\ln\sqrt{\frac{\pi}{\beta a}}
= \frac{1}{2\beta},
\end{split}\end{equation}
where the contribution of $E'$ cancels between numerator and denominator, and we used the trick of Equation~\eqref{eq:z-derivatives} on the one-variable partition function.
Therefore
\begin{equation}\label{eq:equipartition}
\langle a q^2\rangle = \frac{1}{2}\,k_B T
\qquad\text{for every quadratic degree of freedom,}
\end{equation}
independently of the stiffness $a$ and of the nature of $q$.
In particular, for the kinetic energy of a one-dimensional velocity,
\begin{equation}\label{eq:equipartition-kinetic}
\Bigl\langle \frac{m v^2}{2}\Bigr\rangle = \frac{k_B T}{2},
\end{equation}
which is exactly the relation $\langle mv^2\rangle = k_B T$ inserted into the ensemble-averaged Langevin equation in Chapter~2.
The theorem holds classically for any degree of freedom entering the energy quadratically; its failure at low temperature, where quantum mechanics freezes degrees of freedom out, was one of the historical doors to quantum theory --- but everything in this book lives safely in the classical regime.


\section{Stationary states, detailed balance, and the Einstein relation}
\label{app:stationary}

The third circle of ideas the book borrows concerns \emph{how a stochastic process encodes equilibrium}.
Chapter~3 showed that a Fokker--Planck equation is a continuity equation for the probability,
\begin{equation}\label{eq:fp-continuity-app}
\frac{\partial}{\partial t}P(x,t) + \frac{\partial}{\partial x}j(x,t) = 0,
\qquad
j(x,t) = a_1(x)\,P(x,t) - \frac{1}{2}\frac{\partial}{\partial x}\bigl[a_2(x)\,P(x,t)\bigr].
\end{equation}
A \emph{stationary} state has $\partial_t P = 0$, hence a current constant in $x$; if the density decays at the boundaries the constant must vanish, and the stationary state satisfies
\begin{equation}\label{eq:detailed-balance}
j(x) = 0 \qquad \text{everywhere},
\end{equation}
the condition of \emph{detailed balance}: not only is the net probability flow zero globally, it is zero across every point --- every microscopic exchange is individually balanced by its reverse.
This is the defining property of thermodynamic \emph{equilibrium}, as opposed to a mere steady state (which can carry a constant nonzero current, like a wire carrying a constant electric current).

Detailed balance is a first-order equation for the stationary density, solved by one integration.
For constant diffusion, $a_2 = D$, Equation~\eqref{eq:detailed-balance} gives
\begin{equation}\label{eq:stationary-solution}
P_s(x) \;\propto\; \exp\biggl\{\frac{2}{D}\int^x a_1(x')\,dx'\biggr\},
\end{equation}
and if the drift derives from a potential, $a_1(x) = -\mu_m V'(x)$ with a mobility $\mu_m$, the stationary state
\begin{equation}\label{eq:boltzmann-from-fp}
P_s(x) \;\propto\; \exp\Bigl\{-\frac{2\mu_m}{D}\,V(x)\Bigr\}
\end{equation}
has \emph{exactly the Boltzmann form}~\eqref{eq:barometric-app}, with the combination $D/2\mu_m$ playing the role of $k_B T$.
This is the precise sense in which a zero-flux stationary state of a diffusion ``looks thermal'', whatever the microscopic origin of the noise --- the observation behind the effective \emph{temperature of income} $T_r = B_0/A_0$ in the wealth distribution of Chapter~8.

\medskip

Requiring consistency between this stationary state and true thermal equilibrium produces one of the deepest results of the subject.
Take the Ornstein--Uhlenbeck process for the velocity of a Brownian particle (Chapter~4),
\begin{equation}\label{eq:ou-app}
dv(t) = -\kappa\, v\, dt + \sqrt{D}\, dW(t),
\qquad
\kappa = \frac{6\pi\eta R}{m},
\end{equation}
whose stationary distribution is Gaussian with variance $D/2\kappa$.
But the stationary distribution of a velocity in a bath at temperature $T$ \emph{must} be Maxwell--Boltzmann, with variance $k_B T/m$ by Equation~\eqref{eq:mb-variance}.
Equating the two,
\begin{equation}\label{eq:einstein-relation}
\frac{D}{2\kappa} = \frac{k_B T}{m}
\qquad\Longrightarrow\qquad
D = \frac{2\kappa\, k_B T}{m},
\end{equation}
the \emph{Einstein relation}: the strength $D$ of the fluctuations and the friction $\kappa$ that dissipates them are not independent, but locked together by the temperature.
This is the simplest instance of the \emph{fluctuation--dissipation} theorem --- noise and damping are two faces of the same microscopic collisions --- and it is the same physics as the Einstein--Smoluchowski coefficient $D_{E.S.} = k_B T/6\pi\eta R$ of Chapter~2, there derived by balancing drift and diffusion currents, here by matching a stochastic process to the canonical ensemble.
That a nineteenth-century ensemble and a twentieth-century stochastic differential equation meet on the same formula is, in miniature, the whole method of this book.
